\noindent\textbf{REMARKS. ON NETWORK ADJUSTMENT.}\hfill327

\begin{center}\textbf{REMARKS.}\end{center}

\par\noindent\hspace*{\parindent}Only a number of loose, unnumbered leaves survive from the network adjustment.
A complete reconstruction in every detail on the basis of these leaves
--- if it is possible at all --- would entail lengthy calculations.

\par\noindent\hspace*{\parindent}Art.~[2], the normal equations of Art.~[4], the side equations in Art.~[5],
and the table under [9], to which only the column headings have been added,
are taken directly from these leaves. Art.~[3], the tables in Arts.~[7] and [8],
and the normal equations of Art.~[6] have been compiled from them. Note [1]
is taken from an observation and computation notebook for the Hanoverian
arc measurement dating from 1825.

\par\noindent\hspace*{\parindent}\textsc{Gauss}'s adjustment procedure is as follows.
After setting up the normal equations on pp.~303 and 304 from the angle
equations in accordance with the method of least squares, \textsc{Gauss}
solved them (by an indirect method, see p.~325) and calculated from them
the corrections that satisfy the triangle angle equations \textls[120]{alone}.
With the angle values after this first correction, the 12 side equations
were set up in the manner peculiar to \textsc{Gauss}, pp.~304---311.
By means of the modified side equation, \textsc{Gauss} ensures that whenever
it is unconnected to another, the adjustment affects only the angle sums
of the adjacent triangles, leaving those of the interior triangles of the
figure to which the side equation refers unaffected. See the letter to
\textsc{Gerling} of 19~January 1840, p.~253. From the transformed side
equations the 12 normal equations on p.~312 were derived. Solving them
yields new values for the corrections to the directions that satisfy the
side equations \textls[120]{alone}. After this first adjustment, the newly
arising triangle misclosures are adjusted again. The corrections obtained
are used to set up the side equations and their associated normal equations
anew. These now differ from their predecessors only in the constant terms.
Solving them again gives new values of the corrections to the directions.
In this way \textsc{Gauss} repeated the adjustment twice more.
As shown in Supplementum theor.~comb.~observ.~Art.~19, this procedure
leads to the same values as adjusting all angle and side equations at once.
According to \textsc{Gauss}, the calculation yielded stable results so quickly
owing to his choice of form for the side equations. See p.~251.

\par\noindent\hspace*{\parindent}If one substituted the corrections in Art.~[7] into all 51 angle equations,
then, as can be seen from the tabulation in Art.~[8], 11 equations would
have an error of $0$, 18 equations an error of $0^{\prime\prime}{,}001$,
18 equations an error of $0^{\prime\prime}{,}002$, and 4 equations an error
of $0^{\prime\prime}{,}003$.

\par\noindent\hspace*{\parindent}\textsc{Gauss} applied further corrections to the directions from points 28
through 33, and the origin of these corrections is unclear. They probably
result from a subsequent change to the two angles observed at Langwarden
by $0^{\prime\prime}{,}128$ and $0^{\prime\prime}{,}015$. In any event, the
corrections given under [7], which were used in deriving the following
tables, are those obtained from his adjustment.

\par\noindent\hspace*{\parindent}From the 51 triangle misclosures $w_i$ one obtains for the mean error $m$
of a direction, using the approximate formula
$\displaystyle m^2=\frac{1}{6\mathbin{.}51}\sum_{1}^{51}w_i^2$,
\[
m=\pm\,0^{\prime\prime}{,}643.
\]

\clearpage
\noindent328\hfill\textbf{REMARKS. ON NETWORK ADJUSTMENT.}
\par\medskip

\par\noindent\hspace*{\parindent}In connecting the point Hohenhorn to the triangulation network on p.~316
by means of the figure Hohenhorn.~21.~15.~17.~18, the side ratios and angles
of the arc measurement were held fixed. The triangle side Hamburg--Hohenhorn
served as the baseline for the Hanoverian arc measurement and land survey.
See Bestimmung des Breitenunterschiedes etc., p.~48, and the letter to
\textsc{Schumacher} of 22~December 1827, p.~281.

\par\noindent\hspace*{\parindent}The following should be noted about the table in Art.~[9].
Let the representation in the plane of a triangle side $i.k$ make the
angles $\psi_{i.k}$ and $\psi_{k.i}$ with the straight line through its
starting and ending points. Let $T$ denote the astronomical azimuth, $t$
the azimuth on the spheroid, $\Theta$ the azimuth in the plane, and
$\mathrm{c}$ the meridian convergence. Then (see p.~201)
\[
\begin{array}{r@{\qquad}c}
1)&t_{i.k}=T_{i.k}+c_i\\[4pt]
2)&\Theta_{i.k}=t_{i.k}-\psi_{i.k},\qquad
\Theta_{k.i}=\Theta_{i.k}\pm180^\circ,\qquad
t_{k.i}=\Theta_{k.i}+\psi_{k.i}.
\end{array}
\]

\par\noindent\hspace*{\parindent}The calculation of $\psi$ follows the formula in Art.~15, p.~159.
See also pp.~216/217. This formula, like the one for reducing a spheroidal
triangle side to the plane, requires approximate coordinates of the
triangle vertices, which must be obtained by a preliminary calculation
of low precision.

\par\noindent\hspace*{\parindent}Thus if the astronomical azimuth of \textls[120]{one} triangle side
in the spheroidal triangulation network is known, one can calculate $t$
and $\Theta$ for all sides by successive application of equations (2)
and with the aid of the angles of the network. Now for all points on the
principal meridian $\mathrm{c}=0$. Hence if the triangle side $1\,.\,2$,
whose azimuth has been observed, starts from a point on that meridian, then
\[
\begin{array}{r@{\qquad}c}
3)&t_{1.2}=T_{1.2}.
\end{array}
\]

\par\noindent\hspace*{\parindent}The side Göttingen, Theodolithplatz 1823---Nördliches Meridianzeichen
orients the Hanoverian system of triangles. According to Art.~[9]
its azimuth is $T_{1.2}=t_{1.2}=5^{\prime\prime}{,}473$.
In a letter to \textsc{Gerling} of 26~December 1823 (see also \textsc{Gerling},
Beiträge zur Geographie Kurhessens etc., p.~69), \textsc{Gauss} gives
$5^{\prime\prime}{,}471$ for it. The coordinates of the 1823 theodolite
station at the Göttingen Observatory, $x=-5{,}507\,\mathrm{m}$, $y=0$, and
those of the northern meridian mark, $x=-5019{,}756\,\mathrm{m}$,
$y=-0{,}133\,\mathrm{m}$, Volume~IV, pp.~416/417, likewise give the azimuth
$=5^{\prime\prime}{,}471$.

\par\noindent\hspace*{\parindent}Once the plane azimuths have been computed, as in the table under [9],
one can use them to assemble the angles of the plane triangles and then,
if the linear length of one triangle side is known, calculate the lengths
of all the other triangle sides. One therefore need only \textls[120]{once}
reduce a spheroidal triangle side to a plane one. This conversion is made
by the formula given on p.~215. Applying it to all spheroidal triangle
sides likewise gives the plane triangle sides.

\par\noindent\hspace*{\parindent}Once the sides of the plane triangles and their azimuths are known, one
can readily use them to calculate the plane rectangular coordinates of
the triangle vertices in succession.

\par\noindent\hspace*{\parindent}The map on p.~299 of the triangulation network subjected to adjustment is
a copy of a drawing that \textsc{Gauss} submitted to the Hanoverian cabinet
ministry. The directions Wilsede-Brüttendorf and Timpenberg-Hamburg were
omitted in the original by oversight. The original does, however, also show
the connection of the points Hohenhorn and Lauenburg, and that of the points
Wangeroog and Neuwerk.

\par\noindent\hspace*{\parindent}For the adjustment of the triangles observed in the first years of the
arc measurement, see the letters below to \textsc{Bessel} of 5~November 1823
and to \textsc{Bohnenberger} of 16~November 1823.
\hfill\textsc{Krüger.}

\begin{center}\rule{.24\textwidth}{.4pt}\end{center}
